\par\noindent\begin{minipage}{\linewidth}
% !TeX root = ../教材题库_全部条目_转写工作稿.tex
% Source: XJ-B1-C01-HW-P023-Y1381
% 第1章 / 1.2 常用逻辑用语 / 教材 p.23、24 / PDF 第30、31页
% 已对照原页转写题干及公式；待教师复核；未制作答案
% 原题及相邻原解截图（核对用）：../题目原图/XJ-B1-C01-HW-P023-Y1381-01.webp, ../题目原图/XJ-B1-C01-HW-P023-Y1381-02.webp
% 本文件只在转写完成后作为题干使用；OCR 条目尚不是可靠题干。
\begin{bankitem}{习题1.2 第4题}{第1章 / 1.2 常用逻辑用语 / 教材 p.23、24 / PDF 第30、31页}{XJ-B1-C01-HW-P023-Y1381}{已对照原页转写题干及公式；待教师复核；未制作答案}
判断下列结论是否成立，并说明理由。
\subq{1}$p$：$a>b$，$q$：$ac^2>bc^2$，$p$ 是 $q$ 的充分条件；
\subq{2}$p$：$x^2\ne1$，$q$：$x\ne1$，$p$ 是 $q$ 的必要条件；
\subq{3}$p$：四边形为正方形，$q$：四边形为矩形，$p$ 是 $q$ 的充分而不必要条件；
\subq{4}$p$：$a>b$，$q$：$a+c>b+c$，$p$ 是 $q$ 的充要条件；
\subq{5}$p$：$x\in A\cap B$，$q$：$x\in A$，$p$ 是 $q$ 的充要条件。
\end{bankitem}
\end{minipage}\par\medskip
